\documentclass[10pt]{article}
\providecommand{\FIGDIR}{fig}
\newcommand{\DIGESTDATE}{10 November 2026}
\input{preamble}

\begin{document}

% ==================================================================== MASTHEAD
\thispagestyle{empty}
\begin{tikzpicture}[remember picture, overlay]
  \fill[Deep] (current page.north west) rectangle
      ($(current page.north east)+(0,-67mm)$);
  \fill[Amber] ($(current page.north west)+(0,-67mm)$) rectangle
      ($(current page.north east)+(0,-68.4mm)$);
  \foreach \x/\y/\r in {22/12/0.9, 41/26/1.5, 63/9/0.7, 88/31/1.1, 112/17/1.9,
                        138/29/0.8, 159/13/1.3, 178/54/1.0, 196/21/1.6,
                        30/44/0.6, 74/57/1.2, 125/48/0.9, 168/6/0.8, 191/59/0.7,
                        52/18/1.0, 100/60/0.6, 148/20/1.1, 66/49/1.4, 205/45/0.8}
    \fill[Sky, opacity=0.30] ($(current page.north west)+(\x mm,-\y mm)$) circle (\r mm);
\end{tikzpicture}

\vspace*{4mm}
{\sffamily\fontsize{7.6}{9}\selectfont\color{Sky}%
 AUTOMATED DAILY LITERATURE DIGEST\;\textbullet\;PhD RESEARCH WATCH\par}
\vspace{2.2mm}
{\sffamily\bfseries\fontsize{27}{29}\selectfont\color{white}%
 Particulate Matter\par}
{\sffamily\fontsize{27}{29}\selectfont\color{Sky}%
 Monitoring \& Health Effects\par}

\vspace{5.0mm}
{\sffamily\fontsize{8.4}{10}\selectfont\color{white}%
 Issue 2026-11-10\;\;\textcolor{Amber}{\textbar}\;\;Entry window 10 November 2026 (full day)\;\;%
 \textcolor{Amber}{\textbar}\;\;26 records screened in scope\par}
\vspace{18mm}

\noindent
\begin{minipage}[t]{0.190\linewidth}\metric{26}{RECORDS IN SCOPE\\(66 EXCLUDED)}{Deep}\end{minipage}\hfill
\begin{minipage}[t]{0.190\linewidth}\metric{8}{OF 10 SUBTOPIC\\BANDS POPULATED}{Teal}\end{minipage}\hfill
\begin{minipage}[t]{0.190\linewidth}\metric{9}{EFFECT ESTIMATES\\WITH A CI}{Sky}\end{minipage}\hfill
\begin{minipage}[t]{0.190\linewidth}\metric{3}{TIER-A\\RECORDS}{Amber}\end{minipage}\hfill
\begin{minipage}[t]{0.190\linewidth}\metric{4}{RECORDS CARRIED ON\\METADATA ALONE}{Clay}\end{minipage}

\vspace{6mm}

\section*{\sffamily\bfseries\color{Deep}Signal of the day}
\vspace{-1.5mm}{\color{Amber}\rule{\linewidth}{1.1pt}}\vspace{2.5mm}

\begin{keybox}
{\sffamily\bfseries\small\color{Deep}In 55,223 hemodialysis patients across California, Oregon and Washington, a smoke plume over the clinic raised same-day mortality (RR 1.19, 1.11--1.27), but plume plus an extreme heat event gave RR 1.92 (1.71--2.17) against 1.02 (0.95--1.10) for plume alone; per 10\,$\mu$g/m$^3$ of wildfire PM$_{2.5}$ the mortality RR was 1.03 (1.02--1.06).}\\[1.4mm]
\footnotesize
University of Maryland, \textit{Environmental Epidemiology}; Fresenius Kidney Care records 2010--2018, time-stratified case-crossover, conditional quasi-Poisson, exposures assigned to the dialysis facility.
\textbf{Where it is weak:} exposure is clinic-sited, not residential, and wildfire PM$_{2.5}$ is a modelled split of total PM$_{2.5}$; the plume-alone null with a large joint effect means the headline plume signal is mostly heat-coincident, and the abstract gives no wildfire-PM$_{2.5}$ $\times$ heat interaction on the concentration scale.
\textbf{Why it matters here:} dialysis patients are at the same address three times a week. A low-cost monitor on each clinic is an exact exposure point for one of the most smoke-vulnerable populations, and it would replace the modelled split with a measured one.\par
\end{keybox}

\vspace{3mm}

\section*{\sffamily\bfseries\color{Deep}What changed today}
\vspace{-1.5mm}{\color{Amber}\rule{\linewidth}{1.1pt}}\vspace{2.5mm}

\footnotesize
\begin{itemize}

\item \textbf{Nitrate leads a component mixture.} Pearl River Cohort, 182,009 participants: each quartile of a six-component PM$_{2.5}$ mixture index raised respiratory mortality (HR 1.29, 1.23--1.34), nitrate weight 0.48, then organic matter and chloride. Single-component HRs of 1.42--2.75 come without increments in the abstract. The same day, a Peking group's global HONO reconstruction names HONO change as a major driver of particulate nitrate trends in Beijing and Delhi.

\item \textbf{An AI air-quality model beats CAMS at stations but goes negative.} Independent evaluation of the released Aurora 0.4$^\circ$ checkpoint, 96 initialisations in 2024: MSE skill over persistence 38.1--42.6\,\% for PM$_{2.5}$ and 49.1--52.3\,\% for PM$_{10}$, station MAE below CAMS in all six region--pollutant pairs, yet 6.27\,\% of PM$_{2.5}$ and 7.84\,\% of PM$_{10}$ grid values negative at +12\,h and no forecast distribution. Background stations only (preprint).

\item \textbf{Two cohort nulls and one unscaled positive.} Brigham MS cohort (3,040, mean PM$_{2.5}$ 5.8\,$\mu$g/m$^3$): PIRA HR 1.01 (0.96--1.06), EDSS\,$\geq$\,6 HR 1.03 (0.94--1.13). UK Biobank substance use disorder: PM$_{2.5}$ HR 1.111 (1.081--1.142), PM$_{10}$ 1.070 (1.054--1.086), increment not stated in either abstract, so neither can be pooled.

\item \textbf{Sensing side is all modelling.} Four sensing records today --- Aurora evaluation, size-resolved UFP forecasting in Shenzhen (best NRMSE 60.2\,\% at 10\,nm vs 13.4--14.8\,\% at 19--100\,nm), desert AOD forecasting, and a denoising neural-process interpolator that cut RMSE up to 30\,\% on PM$_{2.5}$ from real-world air-quality sensors (preprint, network unnamed). No co-location or calibration record.

\item \textbf{Exposure basis (x1).} One record rests on a low-cost sensor network (the neural-process interpolation benchmark, sensing side) and none on a personal monitor; two rest on regulatory monitors (urolithiasis case-crossover, German ED protocol). On the health side, of 11 records with a recoverable basis 36\,\% are measured (mostly laboratory exposures) and 64\,\% modelled or satellite-fused; \emph{every one of today's nine effect estimates rests on modelled exposure} (wildfire split, component surfaces, LUR, satellite zip-code means).

\end{itemize}

\vspace{2mm}

\begin{keybox}
{\sffamily\bfseries\footnotesize\color{Deep}Window continuity}\\[1.2mm]
{\fontsize{7.2}{8.8}\selectfont
The 9 November issue closed with \texttt{last\_entry\_date} at 9 November. This issue was built on the \textbf{12 November 23:38 UTC (17:38 CST) run}. Before 22:00 CST the newest buildable day is yesterday, so this run builds \textbf{10 and 11 November as two issues}; \textbf{10 November is covered in full}. Tuesday: no rollup is owed; \textbf{W46 (8--14 November) falls on Saturday 14 November}. A cross-day pass over the 10--11 November candidates assigned each record to its earliest day and removed 12 records already seen or rejected. No \textsc{backfill} records today. OpenAlex, alphaXiv, bioRxiv/medRxiv and Wiley Scholar Gateway connectors were not available on this session; the OpenAlex harvester leg returned HTTP 429.\par}
\end{keybox}

\vspace{3mm}
\par\ifdim\dimexpr\pagegoal-\pagetotal\relax<18\baselineskip\newpage\fi
\section*{\sffamily\bfseries\color{Deep}Corpus analytics}
\vspace{-1.5mm}{\color{Amber}\rule{\linewidth}{1.1pt}}\vspace{2.5mm}

\noindent\begin{minipage}{\linewidth}\includegraphics[width=\linewidth]{\FIGDIR/f1_subtopics.png}
\figcap{Subtopic assignment of the 26 records across eight of ten bands: exposure assessment 7, burden 5, sensing 4, other clinical 3, neuro 3, respiratory 2, and one each for reproductive and mechanistic. Cardiovascular and occupational are empty today.}\end{minipage}

\vspace{2mm}
\noindent\begin{minipage}{\linewidth}\includegraphics[width=\linewidth]{\FIGDIR/f2_design_endpoint.png}
\figcap{Left --- architecture: modelling 7 (all four sensing records among them), cohort 5, metadata-only 4, and two each for acute designs, measurement campaigns, experiments and reviews, plus one ecological and one cross-sectional record. Right --- every recoverable endpoint is a singleton; the bar of 4 is the metadata-only records, and eight endpoints are measurement or forecast metrics rather than clinical outcomes.}\end{minipage}

\vspace{2mm}
\noindent\begin{minipage}{\linewidth}\includegraphics[width=\linewidth]{\FIGDIR/f3_metric_geography.png}
\figcap{Left --- particle metric: PM$_{2.5}$ alone 13, PM$_{10}$ alone 4, PM$_{2.5}$ and PM$_{10}$ jointly 4, unspeciated PM or emissions 3, and one each for ultrafine number and optical (AOD). Right --- geography by region: global or multi-region 10 (including the two multi-country Asian records and two with no stated location), China 7, North America 4, Europe 3 (Italy, Germany, UK), and one each for East Asia ex-China (South Korea) and South Asia (India).}\end{minipage}

\vspace{2mm}
\noindent\begin{minipage}{\linewidth}\includegraphics[width=\linewidth]{\FIGDIR/f6_lifecourse.png}
\figcap{Life-course windows. Preconception is the mouse ovarian study; early childhood the GBD analysis, where PM$_{2.5}$-attributable mortality is highest under age 5. Working age carries seven records (UK Biobank, MS cohort, Pearl River, hemodialysis, urolithiasis, arthritis cohort, ED protocol); older adults four (2050 burden projection, Pearl River, arthritis cohort, hemodialysis). No adolescent record.}\end{minipage}

\vspace{2mm}
\noindent\begin{minipage}{\linewidth}\includegraphics[width=\linewidth]{\FIGDIR/f4_heatmap.png}
\figcap{Subtopic $\times$ study architecture. Sensing is entirely modelling (4); exposure assessment splits across modelling (2), metadata-only (2), measurement campaigns (2) and the GIS exposome study (1); burden mixes modelling, the ED protocol (cohort), two metadata-only records and the GBD greenspace analysis. Both acute designs are other-clinical (hemodialysis, urolithiasis).}\end{minipage}

\vspace{2mm}

\IfFileExists{\FIGDIR/extra_figs.tex}{\input{\FIGDIR/extra_figs.tex}}{}



\noindent\begin{minipage}{\linewidth}\includegraphics[width=\linewidth]{\FIGDIR/f5_forest.png}
\figcap{Ratio scale, 9 estimates from 4 sources: Song hemodialysis wildfire (4: two per 10\,$\mu$g/m$^3$ wildfire PM$_{2.5}$, two categorical plume contrasts), Huang Pearl River (1, per quartile of the component mixture), Liu UK Biobank (2, increment not stated) and Lebiedz MS cohort (2, increment not stated). Four of nine estimates carry no stated increment, and the two largest are categorical contrasts; endpoints differ across sources, so no pooling and no I$^2$.}\end{minipage}

\vspace{4mm}

\par\vskip 0pt plus 24\baselineskip\penalty-200\vskip 0pt plus -24\baselineskip
\section*{\sffamily\bfseries\color{Deep}Paper-level digest}
\vspace{-1.5mm}{\color{Amber}\rule{\linewidth}{1.1pt}}\vspace{2.5mm}

{\fontsize{7.4}{9}\selectfont\color{Slate}Ordered by cluster. Each entry gives the
methodological core and the finding that matters, not an abstract paraphrase. Records
marked \textit{metadata only} carry no recoverable abstract and assert no finding.\par}

\band{Neuro / mental health\hfill 3 records}{Deep}

\paperentry{Liu Xi'an UK Biobank PM$_{2.5}$ PM$_{10}$ PMcoarse NO$_2$ NOx greenness incident substance use disorder Cox 10.9 y}
{Various ambient air pollutants, residential greenness, and risk of substance use disorders: a population-based cohort study.}
{BMC Psychiatry}
{247,138 UK Biobank participants free of substance use disorder at baseline; 4,982 incident cases over a median 10.9 years. All five pollutants were positively associated; PM$_{2.5}$ strongest (HR 1.111, 1.081-1.142), then PM$_{10}$ (1.070, 1.054-1.086); the increment is not stated in the abstract, so the size of these HRs cannot be interpreted. Green space (HR 0.997, 0.996-0.998) was inversely associated and domestic garden attenuated the pollutant HRs in the high-exposure group. Effects were largest in older men and ever-smokers, the strata where residual confounding by deprivation and smoking intensity is worst; baseline LUR surfaces stand in for decade-long exposure. Relevance: low - an LUR cohort result without an increment adds a hazard ratio that cannot be pooled.}
{10.1186/s12888-026-08588-8}

\paperentry{Lebiedz Chitnis Brigham MS Center PM$_{2.5}$ satellite zip-code PIRA EDSS 6 Cox null northeastern US 2017-2023}
{Particulate Matter is Not Associated With Disease Worsening in Patients With Multiple Sclerosis: Retrospective Results From a Clinic-Based Cohort in the Northeastern United States.}
{Brain Behav}
{3,040 people with MS at the Brigham MS Center, mean age 49.5, median EDSS 2.7, mean baseline PM$_{2.5}$ 5.8 $\mu$g\,m$^{-3}$. Adjusted HRs were null for PIRA (baseline 1.01, 0.96-1.06; time-varying 0.98, 0.93-1.04) and for EDSS >= 6 (baseline 1.03, 0.94-1.13; time-varying 1.01, 0.92-1.10). The increment behind the HRs is not stated. Exposure contrast is narrow at New England levels, zip-code assignment blurs it further, and clinic-based retrospective ascertainment of PIRA is noisy, so the null mainly bounds large effects at low concentrations. Relevance: a well-powered null at \textasciitilde{}6 $\mu$g\,m$^{-3}$; resolving within-zip contrast would need denser measurement than satellite surfaces give.}
{10.1002/brb3.71878}

\paperentry{Guilin critical review neurovascular glial recovery after particulate exposure reduction filtration diesel BBB microglia}
{Asynchronous neurovascular and glial responses after particulate air pollution exposure declines.}
{Neurotoxicology}
{Critical review of experimental post-exposure studies and human exposure-reduction studies, separating cellular state, direct function and attribution. Early blood-brain barrier leakage after diesel exhaust is documented but later normalisation of the same function is not; some inflammatory markers attenuate after a switch to filtered air while other abnormalities persist for weeks or months. Human studies give both favourable and null cognitive outcomes after exposure reduction, without cell-specific recovery. The review argues for asynchronous responses rather than a common recovery sequence and sets criteria (repeated measurement of one task) for claiming recovery. Relevance: indirect - filtration-intervention trials need verified exposure drops, which indoor sensors provide.}
{10.1016/j.neuro.2026.103598}

\par\vskip 0pt plus 10\baselineskip\penalty-200\vskip 0pt plus -10\baselineskip
\band{Respiratory \& allergic\hfill 2 records}{Sky}

\paperentry{Huang Wei Pearl River Cohort PM$_{2.5}$ components NH4 NO$_3$ BC OM SO4 Cl respiratory mortality quantile g-computation South China}
{Long-Term Exposure to PM$_{2.5}$ Components and Respiratory Disease Mortality in South China: The Pearl River Cohort Study}
{Environ Sci Technol Lett}
{182,009 Pearl River Cohort participants; marginal structural Cox models for single components and quantile g-computation for the mixture. All six components raised respiratory mortality, single-component HRs 1.42-2.75 (increments not given in the abstract). Each quartile increase in the mixture index gave HR 1.29 (1.23-1.34), with nitrate carrying the largest weight (0.48), then organic matter and chloride. Older adults, non-smokers and non-exercisers were most susceptible (HRs 1.25-2.02). Component surfaces are model products whose errors are correlated across species, so the weight ranking partly reflects which component the model resolves best; HRs of 1.4-2.8 per component are implausibly large without the increment. Relevance: nitrate is the semi-volatile, hygroscopic fraction behind most low-cost-versus-reference disagreement in humid seasons.}
{10.1021/acs.estlett.6c00573}

\paperentry{Peking University review ozone PM$_{2.5}$ heat coexposure triple exposure airway inflammation mechanisms immune-neuro-microbiota}
{Mechanisms of Airway Inflammatory Diseases Under Ozone Single Exposure, Coexposure, and Triple Exposure With PM$_{2.5}$ and Heat: A Review.}
{J Appl Toxicol}
{Review of mechanisms linking ozone alone, ozone with PM$_{2.5}$ or heat, and all three to airway inflammation. Ozone injures epithelium oxidatively and activates NF-kB, MAPK and NLRP3; PM$_{2.5}$ is framed as a chemical amplifier giving additive or synergistic oxidative burden, whereas heat acts as a biological sensitiser through TRP-mediated neuroimmune activation and microbiota-metabolism shifts. Triple exposure is argued to overwhelm compensation through higher effective dose, lower tolerance and convergent signalling, but the authors state its health effects are largely unexplored. No quantitative synthesis and no search protocol are described. Relevance: low for sensing; co-located O$_3$, PM$_{2.5}$ and temperature are what multi-pollutant low-cost nodes already log.}
{10.1002/jat.70501}

\par\vskip 0pt plus 10\baselineskip\penalty-200\vskip 0pt plus -10\baselineskip
\band{Reproductive \& developmental\hfill 1 record}{Sage}

\paperentry{Shandong mouse intratracheal PM$_{2.5}$ ovary single-cell RNA-seq leukocyte infiltration granulosa cholesterol biosynthesis STAT1 atresia}
{Single-cell transcriptomic profiling of PM$_{2.5}$ ovarian toxicity: pulmonary particulate instillation associates with leukocyte infiltration and impaired granulosa cholesterol biosynthesis.}
{Toxicology}
{Mouse model of repeated pulmonary PM$_{2.5}$ instillation with ovarian single-cell RNA-seq and in situ validation. Instillation expanded the ovarian immune compartment (B-lineage lymphocytes, neutrophils); resident macrophages took a C1qa-rich scavenger profile without NLRP3 activation. Immune-derived TGF-beta1 and TNF signalling raised JAK-STAT activity in somatic cells and suppressed the granulosa cholesterol-biosynthesis programme; trajectory modelling nominated STAT1 for follicular atresia, with p-STAT1 and TUNEL confirming in tissue. Dose, particle source and group sizes are not in the abstract, and bolus instillation is not inhalation. Relevance: low for sensing; extends the lung-to-ovary immune route behind epidemiological fecundity signals.}
{10.1016/j.tox.2026.154623}

\par\vskip 0pt plus 10\baselineskip\penalty-200\vskip 0pt plus -10\baselineskip
\band{Mechanistic toxicology\hfill 1 record}{Violet}

\paperentry{Korea University PM$_{10}$ brain vascular endothelial cells apoptosis DNA damage MAPK NF-kB NAC zebrafish developmental toxicity}
{Particulate Matter Induces Apoptotic Death of Brain Vascular Endothelial Cells in Association with Oxidative Damage, DNA Damage Response, and MAPK/NF-B Signaling.}
{Environ Toxicol Pharmacol}
{PM$_{10}$ reduced brain vascular endothelial cell viability and raised TUNEL labelling, cleaved caspase-3 and PARP, intrinsic and death-receptor apoptotic proteins, DNA-damage-response markers, MAPK/NF-kB components and lipid and protein oxidation. N-acetylcysteine pretreatment improved survival and reduced caspase-substrate cleavage, implicating redox processes. In zebrafish, PM$_{10}$ produced concentration-dependent developmental toxicity and altered brain vascular phenotypes. The authors themselves describe the causal chain as unresolved; doses, particle provenance and whether concentrations approach realistic deposition are not given in the abstract. Relevance: low - mechanism plausibility for the cerebrovascular and neurodevelopmental associations seen in cohorts.}
{10.1016/j.etap.2026.105230}

\par\vskip 0pt plus 10\baselineskip\penalty-200\vskip 0pt plus -10\baselineskip
\band{Exposure assessment \& modelling\hfill 7 records}{Clay}

\paperentry{Anhui WRF-Chem GOCART Shao KOK dust schemes February 2026 East Asian dust storm size-resolved PM$_{10}$ Gobi source attribution}
{Numerical Analysis of an Early-Arriving East Asian Dust Storm in 2026: Evaluation, Size-Resolved Dust Contribution, and Source Attribution}
{Atmosphere}
{WRF-Chem case study of the 18-23 February 2026 dust storm that reached the Yangtze basin, comparing GOCART, Shao2001 and KOK emission schemes. A tuned GOCART run matched observed AOD best (R 0.74, RMSE 0.16). Size bins 5-8 supplied over 90\% of dust mass in surface PM$_{10}$; regional emission peaked above 4.0 Tg on 20 February while area-mean loading peaked at about 90 mg/m2 a day later; the Gobi supplied about two-thirds of the loading. Validation is against AOD, not surface PM$_{10}$ monitors, and one event cannot rank schemes. Relevance: coarse-dominated dust episodes are where PM$_{2.5}$-tuned optical sensors under-read PM$_{10}$ most; event-specific size splits like this are what a dust correction needs.}
{10.3390/atmos17111107}

\paperentry{Zhang Peking deep neural network global surface HONO reconstruction 2005-2021 megacities ozone particulate nitrate drivers}
{Long-term nitrous acid evolution modulates ozone and secondary aerosol variations in global megacities.}
{Nat Commun}
{A deep neural network trained on worldwide field campaigns reconstructs surface HONO; Beijing fell 47.4\% (2013-2021), Delhi rose 10.0\% (2005-2018), and London, New York and Los Angeles declined slightly. HONO change is identified as a major driver of ozone and particulate nitrate variation, outweighing NOx and VOC abatement in Beijing and Delhi, and the framework recommends NOx control in summer and joint NOx, PM$_{2.5}$ and direct-combustion cuts in winter. HONO is sparsely measured and campaign-biased, so the reconstruction and the attribution share one learned model; no uncertainty on the nitrate share is given in the abstract. Relevance: indirect - explains nitrate-driven PM$_{2.5}$ trends that sensor networks record but cannot apportion.}
{10.1038/s41467-026-78349-4}

\paperentry{NIEHS PEGS GIS-ExWAS 88 geospatial exposure indicators 11 phenotypes IPTW LASSO air pollution cardiovascular}
{Applying geographic information systems data linkages for an exposome-wide association study in the Personalized Environment and Genes Study}
{J Expo Sci Environ Epidemiol}
{7,531 Personalized Environment and Genes Study participants' residential histories linked by open-source GIS to 88 exposure indicators; cross-sectional associations with 11 phenotypes by adjusted logistic regression, IPTW and LASSO. The GIS-ExWAS recovered associations complementary to the cohort's earlier questionnaire ExWAS, including air pollution with cardiovascular disease, plus new suggestive associations. No effect sizes are given; the contribution is the linkage pipeline. Cross-sectional design and many correlated indicators make every hit a hypothesis, and the abstract does not say which air-pollution surfaces were linked. Relevance: indirect - a standard linkage layer is where sensor-network surfaces would have to plug in.}
{10.1038/s41370-026-00986-0}

\paperentry{Chi NYCU PM-bound PCDD/F PCB PCN TEQ Ulaanbaatar Beijing Taipei PMF source apportionment cancer risk dust events}
{Spatiotemporal distribution, source apportionment and cancer risk of PM-bound PCDD/Fs, PCBs and PCNs in three East Asian cities with contrasting emissions: Ulaanbaatar, Beijing and Taipei area.}
{Environ Pollut}
{61 TSP samples in Ulaanbaatar, 24 in Beijing (22 PM$_{2.5}$) and 26 PM$_{2.5}$ in the Taipei area. PM$_{2.5}$-bound PCDD/F TEQ averaged 107 fg/m3 in Beijing against 11.5 in Taipei; Ulaanbaatar TSP reached 128. Beijing cold-season PCDD/F TEQ rose 19-fold, and during dust events Beijing TSP TEQ (76.7) exceeded Ulaanbaatar's (48.9). PMF attributed PCDD/Fs to open burning and traffic (Ulaanbaatar), waste incineration and dust (Beijing) and boilers plus incineration and traffic (Taipei). PCDD/Fs drove over 98\% of carcinogenic risk, which exceeded 1e-6 in Ulaanbaatar and Beijing, but non-inhalation pathways carried 65-90\% of it. Sampling periods and size cuts differ by city, limiting comparison. Relevance: indirect - mass-based sensors cannot see the toxic-equivalent fraction that drives risk here.}
{10.1016/j.envpol.2026.129475}

\paperentry{Peanusaha Auvermann Texas cattle feedlot evening dust peak Hot-Dry-Windy index water balance PM$_{10}$ surface moisture}
{A Meteorological Indicator for Particulate Matter Emissions: Adapting the Hot-Dry-Windy Index to Predict Feedlot Evening Dust Peaks}
{Atmos Pollut Res}
{Partial abstract (Crossref). Feedlot evening dust peaks have been documented for decades but their meteorological drivers lack long-term study, and direct surface-moisture measurement is labour-intensive. The authors adapt the Hot-Dry-Windy index and couple it with precipitation in an empirical water-balance model to estimate pen-surface moisture from routine meteorology. Under the driest surface conditions the index explained 43\% of the variance in evening PM$_{10}$ (R2 0.433, r 0.658). The deposited text ends there, so the site, record length, PM$_{10}$ instrument and performance under wetter conditions are unknown. Relevance: a cheap meteorological predictor of agricultural dust episodes, testable against low-cost PM$_{10}$ sensors at feedlot fences.}
{10.1016/j.apr.2026.103257}

\paperentry{Reatini Stafoggia Bayesian hierarchical vs random forest spatiotemporal PM$_{10}$ PM$_{2.5}$ NO$_2$ prediction Italy}
{Comparison of Bayesian hierarchical modelling and Random Forest algorithm for spatiotemporal prediction of PM$_{10}$, PM$_{2.5}$ and NO$_2$ across Italy}
{Atmos Environ}
{\textsc{metadata only.} Metadata only: Elsevier day-0 deposit with no abstract in Crossref at issue time; no finding is asserted. The title states a head-to-head comparison of Bayesian hierarchical modelling and random forest for national spatiotemporal prediction of PM$_{10}$, PM$_{2.5}$ and NO$_2$ across Italy; the author list (Reatini, Nobile, Stafoggia, Stortini) suggests the Italian national exposure-modelling programme. Relevance: direct if the comparison reports uncertainty calibration as well as cross-validated R2 - a Bayesian surface gives the per-cell prediction variance that guides where to add monitors, and random forest gives none without extra machinery. Recheck at the W46 abstract sweep.}
{10.1016/j.atmosenv.2026.122456}

\paperentry{Saurabh Sarangi elemental characterization PM$_{10}$ sources pristine high-altitude western Himalaya}
{Elemental characterization of PM$_{10}$ sources in the pristine high-altitude western Himalayan region}
{Atmos Environ}
{\textsc{metadata only.} Metadata only: Elsevier day-0 deposit with no abstract in Crossref at issue time; no finding is asserted. The title describes elemental characterisation and source identification of PM$_{10}$ at a pristine high-altitude site in the western Himalaya; geography is taken from the title. Sampling method, period, analytical technique and source-apportionment approach are unknown, so whether this is a filter-based XRF or ICP campaign cannot be told. Relevance: background-site composition fixes the dust and long-range transport baseline against which Indo-Gangetic and Himalayan foothill sensor networks are interpreted. Recheck at the W46 abstract sweep.}
{10.1016/j.atmosenv.2026.122492}

\par\vskip 0pt plus 10\baselineskip\penalty-200\vskip 0pt plus -10\baselineskip
\band{Sensing, forecasting \& instrumentation\hfill 4 records}{Amber}

\paperentry{Parthenope independent evaluation Aurora 0.4 degree air pollution foundation model CAMS persistence stations negative PM$_{2.5}$ 2024}
{Independent evaluation of Aurora air pollution forecasts beyond benchmark skill}
{Research Square (preprint)}
{Preprint. 96 initialisations across 2024 of the released Aurora air-pollution checkpoint, scored against CAMS, persistence and hourly background stations in three regions. Relative to persistence, Aurora cut MSE by 38.1-42.6\% for PM$_{2.5}$ and 49.1-52.3\% for PM$_{10}$, and its station MAE beat the benchmark-selected CAMS fields in all six region-pollutant pairs with bootstrap intervals excluding zero. But at +12 h, 6.27\% of PM$_{2.5}$ and 7.84\% of PM$_{10}$ grid values were negative, the model gives no forecast distribution, and its response to perturbed emission inputs is untested against real emission change. Background stations only, so urban hotspots are unscored. Relevance: dense low-cost networks are the obvious test bed for exactly the near-source skill this evaluation could not reach.}
{10.21203/rs.3.rs-11417641/v1}

\paperentry{Shenzhen Mini-WRAS size-resolved ultrafine particle hourly forecasting MLR RF XGBoost LSTM history window 2023-2025}
{Forecasting hourly size-resolved ambient ultrafine particles with size-specific historical windows.}
{Environ Pollut}
{8,982 quality-controlled hourly observations (Sep 2023-Jun 2025) from a Mini-WRAS at an urban Shenzhen site; 270 model-window-target configurations (MLR, RF, XGBoost, LSTM, BiLSTM; 1-24 h histories) scored on the same 533 test hours. Best NRMSE was 60.2\% at 10 nm and 22.1\% at 14 nm against 13.4-14.8\% at 19-100 nm; a 2 h history was best for 10-19 nm and 4 h for 27-100 nm, while the summed UFP target reached 13.7\% at 12 h. Several model classes tied, so architecture mattered less than window. One site, one-step-ahead, and no meteorological or traffic covariates are described. Relevance: aggregate number forecasts hide the poorly predictable nucleation mode, which is also the size range low-cost counters miss.}
{10.1016/j.envpol.2026.129537}

\paperentry{Tohidian Elia denoising neural processes input uncertainty PM$_{2.5}$ spatial interpolation air quality sensors 2026}
{DnNP: Denoising Input Uncertainty in Neural Processes}
{arXiv (preprint)}
{Preprint. Denoising Neural Processes extend the neural-process meta-learner to treat inputs as noisy, trained by variational inference, with a proof that the Bayes-optimal predictor under input noise lies in the model class and that fixed-variance plug-in models carry an irreducible KL gap. On PM$_{2.5}$ spatial interpolation from real-world air-quality sensors the method cut RMSE by up to 30\% and NLL by up to 73\% against standard neural processes. The abstract names neither the network, its sensor type nor the baseline beyond vanilla NPs, so 'real-world sensors' cannot be placed against kriging or calibrated-network practice. Relevance: modelling noisy and mis-located sensor inputs explicitly is exactly the low-cost network problem; worth reading if the PM data are PurpleAir-class.}
{10.48550/arXiv.2611.14424}

\paperentry{Punjab ARIMA LSTM NARX forecasting Deep Blue AOD550 Thar Thal Karakum Kutch Kyzylkum Arabian deserts 2000-2030}
{Forecasting of aerosol optical depth over arid and semi-arid regions using time series and machine learning models.}
{Environ Monit Assess}
{294 monthly Deep Blue AOD550 values (Feb 2000-Jul 2024) for the Thar, Thal, Karakum, Kutch, Kyzylkum and Arabian deserts; ARIMA, LSTM and NARX evaluated on a 44-month test set. NARX was best (test RMSE 0.0180-0.0275) over LSTM (0.0215-0.0350, R2 0.53-0.70) and ARIMA (0.0250-0.0400, R2 0.40-0.60), and was used to project AOD to July 2030. Monthly means of a column optical quantity carry no surface PM information, and a six-year recursive forecast from 250 training points has no credible uncertainty statement. Relevance: marginal - desert dust is the regime where AOD-to-PM$_{2.5}$ conversion and optical low-cost sensors both fail, but this paper does not touch the surface.}
{10.1007/s10661-026-16140-0}

\par\vskip 0pt plus 10\baselineskip\penalty-200\vskip 0pt plus -10\baselineskip
\band{Burden, policy \& mitigation\hfill 5 records}{Coral}

\paperentry{CUHK CESM XGBoost 0.1 degree PM$_{2.5}$ mortality GEMM SSP1 SSP2 SSP3 SSP5 2010 2050 emissions climate land use}
{High-Resolution PM$_{2.5}$ Mortality Projections With Machine Learning Reveal Global Inequities Under Future Climate and Socioeconomic Pathways}
{GeoHealth}
{CESM factorial runs separating emissions, climate and land-use change, bias-corrected and refined to 0.1 degree with XGBoost, fed to GEMM. Present-day burden about 6.13 million deaths in 2010 (UI 5.77-6.48); by 2050 SSP1 removes about 1.1 million deaths and SSP3 adds about 0.5 million. Emissions dominate, climate change offsets part of the emission benefit, and the burden concentrates in people over 60 and in LMICs. The uncertainty interval covers the concentration-response function only; the ML downscaler is trained against present-day observations and is applied to 2050 states it has never seen. Relevance: the downscaler's skill rests on ground data that are thinnest in the LMICs that carry the burden.}
{10.1029/2026gh001829}

\paperentry{Shandong GBD 2021 PM$_{2.5}$ attributable ASMR greenspace access Gini 204 countries 2000-2020 EAPC}
{Inequalities in greenspace access and their associations with all-cause and cause-specific mortality attributable to PM$_{2.5}$: a global analysis from 2000 to 2020.}
{BMC Public Health}
{204 countries, 2000-2020. Global PM$_{2.5}$-attributable all-cause ASMR fell from 63.5 to 55.4 per 100,000 while greenspace access fell from 37.0\% to 28.6\%; greenspace inequality in low-SDI countries rose (Gini 0.20 to 0.24). Higher greenspace access correlated with lower PM$_{2.5}$-attributable mortality for IHD, stroke, COPD and neonatal disorders. Both sides are modelled country aggregates and the attributable burden is itself a function of GBD exposure and baseline rates, so the correlation is ecological and confounded by development. Mortality attributable to PM$_{2.5}$ can fall while concentrations rise if baseline rates fall, so the trend is not an exposure trend. Relevance: low - burden bookkeeping with no measurement content.}
{10.1186/s12889-026-28654-w}

\paperentry{Luebeck Schleswig-Holstein emergency department climate air quality infectious surveillance DLNM register protocol 2014-2024}
{Climate, Air Quality, and Emergency Department Use in Northern Germany: Protocol for a Multicenter Register Study.}
{JMIR Res Protoc}
{Protocol for a four-site retrospective study of emergency department presentations in the Luebeck area, January 2014-December 2024, linking ICD-coded visits to German Weather Service meteorology, state air-quality monitoring and Robert Koch Institute infectious surveillance. Distributed lag nonlinear models within GAMs, stratified by age, sex and socioeconomic status; ethics granted July 2025 and data transfer under way. No results. The analysis plan centres on temperature; the abstract names no particle metric, and a maritime region with few monitors will give little PM contrast. Relevance: low - noted as an upcoming low-concentration ED series.}
{10.2196/92900}

\paperentry{Liu Guanzhong Basin agricultural reactive nitrogen ammonia emissions air quality impact}
{Impacts of agricultural reactive nitrogen emissions on the air quality in the Guanzhong Basin, China}
{Atmos Pollut Res}
{\textsc{metadata only.} Metadata only: Elsevier day-0 deposit with no abstract in Crossref at issue time; no finding is asserted. The title addresses the impact of agricultural reactive nitrogen emissions on air quality in the Guanzhong Basin (Xi'an region); whether PM$_{2.5}$ is the endpoint is not stated in the metadata, and it is carried because agricultural ammonia is the lever on secondary ammonium nitrate in this basin. Relevance: ties to today's nitrate theme (Pearl River components, HONO reconstruction) if the abstract confirms a PM$_{2.5}$ outcome; recheck at the W46 abstract sweep. Model or observation basis is unknown.}
{10.1016/j.apr.2026.103267}

\paperentry{Li Cao plant communities relative reduction efficiency atmospheric microplastics PM$_{2.5}$ urban green space}
{Assessing the relative reduction efficiency of different plant communities to atmospheric microplastics and PM$_{2.5}$}
{Atmos Pollut Res}
{\textsc{metadata only.} Metadata only: Elsevier day-0 deposit with no abstract in Crossref at issue time; no finding is asserted. The title compares how different plant communities reduce atmospheric microplastics and PM$_{2.5}$, presumably by paired in- and out-of-canopy measurement, but neither design, site, season nor instrument is available, and microplastic counting methods are not standardised. Relevance: vegetation-barrier claims rest on short paired campaigns where sensor inter-unit bias is the same size as the effect; worth checking how the reduction was measured, and over how many paired days, once the abstract appears.}
{10.1016/j.apr.2026.103268}

\par\vskip 0pt plus 10\baselineskip\penalty-200\vskip 0pt plus -10\baselineskip
\band{Other clinical endpoints\hfill 3 records}{Slate}

\paperentry{Song Fresenius hemodialysis ESKD wildfire smoke plume PM$_{2.5}$ extreme heat case-crossover mortality West US 2010-2018}
{Exposure to wildfire smoke and risk of mortality and hospitalization among hemodialysis patients in the Western United States: Effect modification by extreme heat events.}
{Environ Epidemiol}
{55,223 adult hemodialysis patients at Fresenius clinics in California, Oregon and Washington, warm seasons 2010-2018; conditional quasi-Poisson. A smoke plume overhead raised same-day mortality (RR 1.19, 1.11-1.27); each 10 $\mu$g\,m$^{-3}$ of wildfire PM$_{2.5}$ raised mortality 3\% (RR 1.03, 1.02-1.06) and hospitalization 1\% (1.01, 1.00-1.03). Plume plus extreme heat gave RR 1.92 (1.71-2.17) against 1.02 (0.95-1.10) for plume alone, so the plume signal is mostly heat-coincident. Exposure is assigned to the clinic, not the residence, and wildfire PM$_{2.5}$ is a modelled split of total PM$_{2.5}$; the interaction contrast is large enough that residual confounding by heat timing remains the main threat. Relevance: a dialysis population attends clinics thrice weekly, which makes clinic-sited low-cost monitors a cheap, exact exposure point.}
{10.1097/ee9.0000000000000541}

\paperentry{Yan Li Sun Yat-sen national cohort PM$_{2.5}$ components OM BC SO4 NO$_3$ NH4 arthritis quantile g-computation middle-aged older China}
{Component-Specific PM$_{2.5}$ Exposure and Arthritis Risk in China: Evidence of Differential Susceptibility in a National Cohort}
{GeoHealth}
{14,935 middle-aged and older adults in a nationally representative Chinese cohort; mixed-effects logistic regression and quantile g-computation. Both short- and long-term PM$_{2.5}$ mass and organic matter, black carbon, sulfate and ammonium were reported to raise arthritis risk, with men and overweight or obese participants most vulnerable. The abstract gives no effect size, CI or increment, and does not say how arthritis was ascertained, so neither magnitude nor outcome specificity can be judged from the record. Component surfaces are correlated model products. Relevance: low until numbers appear; another component study where nitrate does not lead, unlike the Pearl River respiratory result today.}
{10.1029/2025gh001702}

\paperentry{Guangzhou Tianjin Guizhou urolithiasis stone composition NO$_2$ PM$_{10}$ inverse U-shaped case-crossover 128000 stones}
{Ambient nitrogen dioxide and composition-specific risk of urolithiasis: a multicenter time-stratified case-crossover study.}
{BMC Public Health}
{25,374 analysed urinary stones in Guangzhou with validation in Tianjin (99,552) and Guizhou (3,021); conditional logistic regression and GAMs over lags 0-7 and 6- to 12-month averages. The robust signal is NO$_2$, strongest at lags 0 and 2-4 for carbonate apatite and calcium oxalate. PM$_{10}$ showed mainly inverse or U-shaped associations with overall and composition-specific risk, and a PM$_{10}$-carbonate apatite pattern was confirmed in Guizhou; no PM effect sizes are given in the abstract. Stones are dated at analysis, not symptom onset, so short lags are misaligned with formation. Relevance: low - a PM null-to-inverse result in a gas-driven outcome, carried for completeness.}
{10.1186/s12889-026-28561-0}

\clearpage
\section*{\sffamily\bfseries\color{Deep}Corpus register}
\vspace{-1.5mm}{\color{Amber}\rule{\linewidth}{1.1pt}}\vspace{2.5mm}

\input{table_rows}

\vspace{2mm}

\begin{keybox}
{\sffamily\bfseries\footnotesize\color{Deep}Provenance and method}\\[1.2mm]
{\fontsize{7.2}{8.8}\selectfont
Entry window \textbf{10 November 2026}, single day (Tuesday). Harvester legs: \textbf{PubMed} health \textbf{25}, sensing \textbf{3}; \textbf{Europe PMC} \textbf{26}; \textbf{Crossref by ISSN} \textbf{64}; \textbf{arXiv} 1 in-window record (carried); \textbf{OpenAlex} HTTP 429. 119 raw, 111 unique, 17 already seen $\rightarrow$ 94 fresh. PubMed connector (EDAT 10--11 November, three axes) 60 PMIDs, 27 bucketed to 10 November: 6 seen, 1 PMID merged onto a harvested record, 10 appended $\rightarrow$ 104; the cross-day pass removed 12 previously seen or rejected $\rightarrow$ \textbf{92 screened}. Abstracts: efetch filled 15 of 16 PMID records; Crossref recovered a partial abstract for 1 of 5 Elsevier records tried.
OpenAlex, alphaXiv, bioRxiv/medRxiv and Wiley Scholar Gateway connectors were not available on this session; \textbf{Scite} retraction check not run. No Amass or Consensus backfill this run.
\textbf{26 admitted}, \textbf{66 rejected} with logged reasons. \textbf{4 carried on metadata alone}. Affiliations from PubMed XML and Crossref (OpenAlex unavailable). Trial watch: see W46 weekly.
Abstracts remain the copyright of their publishers.\par}
\end{keybox}

\end{document}
